Assembled for the split. The predecessor manuscript's bibliography
(\texttt{Demidont\_KassanjeeBias\_references.bib}) lives on that repository's \texttt{master}
branch and covers only part of this list --- every item in §1 below except
Kassanjee, and almost nothing in §2, entered during the recomputation.

\textbf{All 17 DOIs and 3 PMIDs below were resolved against Crossref or Europe PMC on
29 September 2026}, and the entries added from the reference archive were checked
against the PDFs themselves, not transcribed from the working notes. Two year
discrepancies surfaced and are recorded at the end.

\begin{center}\rule{0.5\linewidth}{0.5pt}\end{center}

\section{1. Estimator and design literature}

\begin{enumerate}
\def\labelenumi{\arabic{enumi}.}
\item
  \textbf{Kassanjee R, McWalter TA, Bärnighausen T, Welte A.} A new general biomarker-based incidence estimator. \emph{Epidemiology} 2012;23(5):721--728. doi:\href{https://doi.org/10.1097/EDE.0b013e3182576c07}{10.1097/EDE.0b013e3182576c07}
\item
  \textbf{Kassanjee R, Pilcher CD, Busch MP, et al.} Viral load criteria and threshold optimization to improve HIV incidence assay characteristics. \emph{AIDS} 2016;30(15):2361--2371. doi:\href{https://doi.org/10.1097/QAD.0000000000001209}{10.1097/QAD.0000000000001209}
  \emph{§3.2: the CEPHIA viral-load optimisation from which Pan's 163 d MDRI derives. Establishes that Pan's basis and our CEPHIA re-estimate share a data lineage and are not independent of each other.}
\item
  \textbf{Gao F, Bannick M.} Statistical considerations for cross-sectional HIV incidence estimation based on recency test. \emph{Statistics in Medicine} 2022;41(8):1446--1461. doi:\href{https://doi.org/10.1002/sim.9296}{10.1002/sim.9296}
  \emph{The framework this paper refines. Assumption C is the point of contact.}
\item
  \textbf{Pan J, Bannick M, Gao F.} Estimating HIV cross-sectional incidence using recency tests from a non-representative sample. \emph{American Journal of Epidemiology} 2026. doi:\href{https://doi.org/10.1093/aje/kwag075}{10.1093/aje/kwag075}
  \emph{Screening-stage selection; the limiting estimation error composed with in §2.7 and recovered in §4.1.}
\item
  \textbf{Wang Q, Duerr A, Gao F.} Addressing population heterogeneity for HIV incidence estimation based on recency test. \emph{Statistics in Medicine} 2025;44(18--19). doi:\href{https://doi.org/10.1002/sim.70216}{10.1002/sim.70216}
  \emph{Covariate transport by reweighting; compared in §S3. Implementation at github.com/qii-wang/HIV-incidence-recency-heterogeneity; the estimator structure used in §S3 was read from that code, not inferred from the paper. Funded by NIH (R01AI177078, R01DA032106, R37AI029168) and Gilead (ISR-US-20-10990), per its acknowledgements --- recorded because the comparison in §S3 is with a method whose development was partly industry-funded, not because it bears on the method's validity.}
\item
  \textbf{Bannick M, Donnell D, Hayes R, et al.} An enhanced cross-sectional HIV incidence estimator that incorporates prior HIV test results. \emph{Statistics in Medicine} 2024;43(17):3125--3139. doi:\href{https://doi.org/10.1002/sim.10112}{10.1002/sim.10112}
  \emph{PT-RITA. Repairs misclassification rather than selection; see Discussion.}
\item
  \textbf{Klock E, Wilson E, Fernandez RE, et al.} Validation of population-level HIV-1 incidence estimation by cross-sectional incidence assays in the HPTN 071 (PopART) trial. \emph{Journal of the International AIDS Society} 2021;24(12). doi:\href{https://doi.org/10.1002/jia2.25830}{10.1002/jia2.25830}
\item
  \textbf{Gao F, Dasgupta S, Pasalar S, et al.} Comparing approaches for estimating counterfactual HIV incidence among populations with high vulnerability to HIV in Lima, Peru: a multi-study comparative analysis. \emph{Journal of the International AIDS Society} 2026;29(9). doi:\href{https://doi.org/10.1002/jia2.70204}{10.1002/jia2.70204}
  \emph{Motivates the counterfactual-placebo application in §1.}
\end{enumerate}

\begin{center}\rule{0.5\linewidth}{0.5pt}\end{center}

\section{2. Empirical parameters}

\begin{enumerate}
\def\labelenumi{\arabic{enumi}.}
\setcounter{enumi}{8}
\item
  \textbf{Duong YT, Kassanjee R, Welte A, et al.} Recalibration of the limiting antigen avidity EIA to determine mean duration of recent infection in divergent HIV-1 subtypes. \emph{PLoS ONE} 2015;10(2):e0114947. doi:\href{https://doi.org/10.1371/journal.pone.0114947}{10.1371/journal.pone.0114947}
  \emph{§3.2: independent MDRI. By binomial regression over \textgreater250 seroconversion panels, 130 d (118--142) at ODn ≤1.5 with PFR 1.6\%; by subtype 129 d (B), 122 d (AE), 109 d (A\&D), 152 d (C). The subtype-C value sits on the Ω\_T* = 151 d of the basis adopted here, from a different panel and method.}
\item
  \textbf{Degenhardt L, Hickman M, Altice FL, et al.} The global epidemiology of injecting drug use, HIV, viral hepatitis and tuberculosis among people who are incarcerated: a multistage systematic review. \emph{International Journal of Drug Policy} 2026;150:105062. doi:\href{https://doi.org/10.1016/j.drugpo.2025.105062}{10.1016/j.drugpo.2025.105062}. PMC13058553
  \emph{§3.5, Route B: IDU prevalence among incarcerated in North America, 13.4\% (95\% CI 10.3--16.8).}
\item
  \textbf{Bradley H, Hall EW, Asher A, et al.} Estimated number of people who inject drugs in the United States. \emph{Clinical Infectious Diseases} 2023;76(1):96--102. doi:\href{https://doi.org/10.1093/cid/ciac543}{10.1093/cid/ciac543}
  \emph{§3.5: denominator of 3.70 M US PWID.}
\item
  \textbf{Feder KA, Sun J, Rudolph JE, et al.} Mortality by cause of death during year 1 of the COVID-19 pandemic in a cohort of older adults from Baltimore, Maryland who have injected drugs. \emph{International Journal of Drug Policy} 2022;109:103842. doi:\href{https://doi.org/10.1016/j.drugpo.2022.103842}{10.1016/j.drugpo.2022.103842}
  \emph{§3.4: ALIVE all-cause mortality, 37.2 and 39.6 per 1,000 person-years.}
\item
  \textbf{Sun J, Mehta SH, Astemborski J, et al.} Mortality among people who inject drugs: a prospective cohort followed over three decades in Baltimore, MD, USA. \emph{Addiction} 2022;117(3):646--655. doi:\href{https://doi.org/10.1111/add.15659}{10.1111/add.15659}
  \emph{§3.4: age-standardised trend, 23 → 45 per 1,000 person-years, 1988--2018.}
\item
  \textbf{Tempalski B, Pouget ER, Cleland CM, et al.} Trends in the population prevalence of people who inject drugs in US metropolitan areas 1992--2007. \emph{PLoS ONE} 2013;8(6):e64789. doi:\href{https://doi.org/10.1371/journal.pone.0064789}{10.1371/journal.pone.0064789}. PMID 23755143
  \emph{§3.8: catchment demographic stability, ρ ≈ 0.}
\item
  \textbf{Gough E, Kempf MC, Graham L, et al.} HIV and hepatitis B and C incidence rates in US correctional populations and high risk groups: a systematic review and meta-analysis. \emph{BMC Public Health} 2010;10:777. doi:\href{https://doi.org/10.1186/1471-2458-10-777}{10.1186/1471-2458-10-777}. PMID 21176146, PMC3016391
  \emph{§3.9: the only directly relevant evidence on η. Overlay band, never a fitted value.}
\item
  \textbf{Handanagic S, Finlayson T, Burnett JC, Broz D, Wejnert C; National HIV Behavioral Surveillance Study Group.} HIV infection and HIV-associated behaviors among persons who inject drugs --- 23 metropolitan statistical areas, United States, 2018. \emph{MMWR Morbidity and Mortality Weekly Report} 2021;70(42):1459--1465. PMID 34673746
  \emph{§3.3: 57\% tested in the preceding 12 months, giving θ = −ln(0.43) = 0.844 y⁻¹. Also §3.5, Route A: past-12-month incarceration 21.0\% and 43.3\%.}
\item
  \textbf{Centers for Disease Control and Prevention.} HIV transmission among male inmates in a state prison system --- Georgia, 1992--2005. \emph{MMWR Morbidity and Mortality Weekly Report} 2006;55(15):421--426. PMID 16628181
  \emph{§3.9: 88 documented seroconversions during incarceration; establishes η\_P \textgreater{} 0.}
\end{enumerate}

\begin{center}\rule{0.5\linewidth}{0.5pt}\end{center}

\section{3. Data sources and software}

\begin{enumerate}
\def\labelenumi{\arabic{enumi}.}
\setcounter{enumi}{17}
\item
  \textbf{CEPHIA public-use dataset.} Consortium for the Evaluation and Performance of HIV Incidence Assays. Zenodo. doi:\href{https://doi.org/10.5281/zenodo.4900634}{10.5281/zenodo.4900634}
  \emph{§3.2: recency function re-estimation. Not redistributed with this work.}
\item
  \textbf{Bannick M, Gao F.} \emph{XSRecency: cross-sectional incidence estimation.} R package, MIT. github.com/mbannick/XSRecency; release 0.2.0, 21 August 2023, commit \texttt{2811002}; \texttt{main} at \texttt{7c1243c5}, 7 July 2025. Companion simulation code: github.com/mbannick/RITA-plus-sims, release \texttt{v081723}. \emph{§3.2 and §4.1. Two contract details were read from the 0.2.0 source rather than assumed. \texttt{get.gamma.params} (\texttt{R/phi-functions.R}) is byte-identical in 0.2.0 and \texttt{main} --- shape = W/(2H−W), rate = 1/(2H−W) --- so our implementation matches either, and \texttt{tests/test\_gamma\_params\_match\_xsrecency} asserts it. \texttt{createRitaCephia} (\texttt{R/get-rita-data.R}) did change: 0.2.0 documents \texttt{ui} as infection duration in days with no conversion, \texttt{main} documents years and divides by 365.25 at line 207. Our procedure works in years and matches \texttt{main}; following the 0.2.0 vignette against \texttt{main} would double-convert.}
\item
  \textbf{Zeng Z.} \emph{Jail Inmates in 2023 -- Statistical Tables.} Bureau of Justice Statistics, US Department of Justice; April 2025. NCJ 309965. \emph{§3.5--3.6: mean 32 d in custody July 2022--June 2023; adult jail incarceration rate by age; all-adult rate 253 per 100,000.}
\item
  \textbf{Kaeble D.} \emph{Time Served in State Prison, 2018.} Bureau of Justice Statistics; March 2021. NCJ 255662. \emph{§3.6: mean time served from initial admission to initial release 2.7 y (median 1.3 y), giving β\_P = 0.37/y.}
\end{enumerate}

\textbf{Bureau of Justice Statistics.} \emph{Prisoners} series. US Department of Justice. \emph{§3.5 Route C only: adult imprisonment rate 453 per 100,000. \textbf{Edition not yet pinned} --- the series is not among the tables committed under \texttt{data/bjs/}, and the 2020 edition reports a COVID-depressed 358 per 100,000, so the figure belongs to a different year. Route C is one of three routes to q and agrees with the BJS-free Route B to 0.04 percentage points, so no conclusion rests on it alone.}

\begin{enumerate}
\def\labelenumi{\arabic{enumi}.}
\setcounter{enumi}{21}
\item
  \textbf{Carson EA.} \emph{Mortality in Local Jails, 2000--2019 -- Statistical Tables.} Bureau of Justice Statistics; December 2021. NCJ 301368. And \emph{Mortality in State and Federal Prisons, 2001--2019 -- Statistical Tables.} Bureau of Justice Statistics; December 2021. NCJ 300953. \emph{§3.7: 167, 330 and 259 deaths per 100,000. Source tables are committed under \texttt{data/bjs/}.}
\item
  \textbf{Maruschak LM.} \emph{HIV in Prisons, 2020 -- Statistical Tables.} Bureau of Justice Statistics; May 2022. NCJ 302601. \emph{Cited only to establish that the incarcerated population is not predominantly PWID, so HIV prevalence cannot substitute for injection prevalence in §3.5. Source tables committed under \texttt{data/bjs/}.}
\item
  \textbf{US Census Bureau.} Population Estimates Program, vintage 2019, county characteristics (\texttt{cc-est2019-alldata}). \emph{§3.10: county-specific 15--64 to 18+ ratios, 0.834--0.908.}
\item
  \textbf{Vera Institute of Justice.} \emph{Incarceration Trends} county-level jail and prison populations. \emph{§3.10: the 2019 anchor and post-2019 jail trajectories. County prison series ends at 2019.}
\item
  \textbf{Kelley CF, Acevedo-Quiñones M, Agwu AL, et al.} Twice-yearly lenacapavir for HIV prevention in men and gender-diverse persons. \emph{New England Journal of Medicine} 2025;392(13):1261--1276. doi:\href{https://doi.org/10.1056/NEJMoa2411858}{10.1056/NEJMoa2411858}
  \emph{PURPOSE 2 (NCT04925752). §3.2, §3.3 and §5.3: the protocol and statistical analysis plan supply the three-month testing exclusion, the adopted assay parameters (MDRI 184 d, FRR 1.5\%, ODn ≤1.5 with VL \textgreater75, T = 2 y), and the plan's own statement that prior testing in the preceding 3--12 months would underestimate the background rate.}
\item
  \textbf{Parkin N, Gao F, Grebe E, et al.} Facilitating next-generation pre-exposure prophylaxis clinical trials using HIV recent infection assays: a consensus statement from the Forum HIV Prevention Trial Design Project. \emph{Clinical Pharmacology and Therapeutics} 2023;114(1):29--40. doi:\href{https://doi.org/10.1002/cpt.2830}{10.1002/cpt.2830}
  \emph{The RAWG recommendations from which the T = 2 y default derives.}
\item
  \textbf{ClinicalTrials.gov.} PURPOSE 4: NCT06101342. Study of lenacapavir and emtricitabine/tenofovir disoproxil fumarate for pre-exposure prophylaxis in people who inject drugs. \emph{§3.10 and §4.7. Phase 2, open-label, multicentre, randomised; pharmacokinetics and safety; eligibility 18 years and older with no upper bound; nine US locations; 181 participants enrolled; start December 2023, actual primary completion July 2026. Registry record consulted 29 September 2026. No efficacy estimate is corrected or commented on; the trial reports none.}
\item
  \textbf{Demidont AC.} Software repository for \emph{Calibration-to-deployment mismatch in HIV prevention trials}. Zenodo 2026. doi:\href{https://doi.org/10.5281/zenodo.20344293}{10.5281/zenodo.20344293}
  \emph{The superseded predecessor analysis. Cited for provenance; its central empirical conclusions are not carried forward. Disclosed under prior publication.}
\end{enumerate}

\begin{center}\rule{0.5\linewidth}{0.5pt}\end{center}

\section{Verification notes}

\textbf{Verified against the supplied PDFs (30 September 2026).} Three checks were run against the
reference archive rather than against metadata:

\begin{itemize}
\tightlist
\item
  \textbf{Degenhardt Table 1, North America row} reads 13·4 (10·3, 16·8) and 246,500 (190,500,
  308,500) --- confirming §3.5 Route B exactly as written. The regional supplement tables carry
  several nearby values for other outcomes (13·0, 13·6) and for the United States alone (13·1);
  these are not the quantity cited and should not be substituted for it.
\item
  \textbf{The Degenhardt supplement records a methodological caveat worth carrying}: for most regions
  a prevalence ratio was pooled only where two or more countries had estimates, but North
  America and Australasia were exceptions where a single country sufficed. The North American
  figure is correspondingly less well supported than its interval suggests.
\item
  \textbf{The recency-basis provenance was traced and one claim withdrawn.} An earlier draft of §3.2
  presented Duong 2015, Pan's 163 d and our CEPHIA re-estimate as three-way corroboration. They
  are not independent: Pan's value traces to Kassanjee 2016, a CEPHIA viral-load optimisation,
  and our re-estimate uses the same consortium data and the same assay-plus-viral-load
  construction. Duong is the only external anchor, and it measures a different algorithm ---
  assay alone, without the viral-load criterion. §3.2 now states the two lineages and why
  Duong's values are systematically shorter.
\item
  \textbf{Handanagic 2021 authorship} was corrected from a corporate CDC attribution to the named
  authors after reading the article's own byline.
\end{itemize}

Two dates in the working notes were wrong in the drafts' favour, and both are now stated as the print year:

{%
\begin{longtable}[]{@{}llll@{}}
\toprule\noalign{}
item & working note & Crossref & used \\
\midrule\noalign{}
\endhead
\bottomrule\noalign{}
\endlastfoot
Bradley & 2023 & issued 2022, \textbf{print 2023}, 76(1):96--102 & 2023 \\
Sun & 2022 & issued 2021, \textbf{print 2022}, 117(3):646--655 & 2022 \\
\end{longtable}
}

Online-first and print years differ for both; the volume numbers confirm the print year.

\textbf{Items 17--24 are grey literature, datasets and software} and carry no DOI. They need the exact report edition, NCJ number where applicable, and an access date before submission. \texttt{docs/audit/computation\_record.md} Appendix C records that the four BJS folders could not be opened locally, and that neither the \emph{Survey of Prison Inmates 2016} (NCJ 252641) nor the 2007--2009 drug-use report contains an injection-route variable --- which is why item 8 supplies that quantity instead.

\textbf{Not cited and deliberately so.} The BJS jail drug-intoxication figure of 184 per 100,000 appears in a page summary and exceeds the all-cause rate, so it is almost certainly a count rather than a rate. It is not used anywhere and should not be cited without checking the underlying table.
